Within the configured candidates and supported excitation tests, the cubic spline is the first candidate for a target-hardware prototype: it retains all three physical inputs with a small coefficient package and is competitive with the full LUT in measured-reference flux error. This is a candidate-selection recommendation, not universal spline superiority or validated application accuracy.

The Gaussian model remains an appropriate timing comparator. Its MATLAB evaluator directly shares exponential values between outputs and derivatives; the PC ranking need not survive a compiled local-spline implementation. The search is limited to deterministic centers, a common width and at most 128 centers. It does not exhaust the RBF design space.

The k representations impose a coordinate-reduction constraint that coefficient compression does not require. A different group split can change their ranking; results are recalculated here rather than transferred from another experiment. Differences from fresh or pre-existing LUTs are not causal error partitions.

No final controller-replacement decision is justified. Missing evidence includes application limits for flux and derivative errors, independent derivative/thermal validation, and target-hardware timing of the actual two-operating-point package with its guard and numeric precision. Temperature was not controlled in the current campaigns, and the secondary artifact cannot close that gap. The sampled sizes, one predetermined partition, restricted coverage and finite k-search budgets limit generalization. No production controller was changed.
